% Unit 049 terjemahan Bahasa Indonesia dari chapter4.tex baris 62--208.
% Sumber: Methods of Algebra, Volume 2, commit
% 9a5803ff77dd3257484cb177f851a73770a59dd3 (CC BY 4.0).
% stable-unit-id: o014.aljabr2.chapter4.triangulated-category-definition
% Cakupan: Bagian 4.1 dari definisi kategori dengan pergeseran sampai catatan
% dualitas; berhenti sebelum lema pada baris 210.

% segment-id: o014.li2.chapter4.triangulated-category-definition.h001
\section{Definisi Kategori Bertriangulasi}\label{sec:triangular-def}
\hypertarget{o014.aljabr2.chapter4.triangulated-category-definition}{ }

% segment-id: o014.li2.chapter4.triangulated-category-definition.p001
Kategori bertriangulasi adalah kategori aditif dengan pergeseran yang
dilengkapi suatu kelas segitiga terbedakan. Terlebih dahulu kita jelaskan apa
yang dimaksud dengan kategori dengan pergeseran.

% segment-id: o014.li2.chapter4.triangulated-category-definition.d001
\begin{definisi}\label{def:category-with-translation}
	\index{kategori!dengan pergeseran (with translation)}
	\index{funktor pergeseran}
	Yang dimaksud dengan \emph{kategori dengan pergeseran} ialah data
	$(\mathcal{D},T)$, dengan $\mathcal{D}$ suatu kategori dan
	$T:\mathcal{D}\to\mathcal{D}$ suatu ekuivalensi.%
	\footnote{Banyak pustaka mensyaratkan bahwa $T$ merupakan automorfisme;
	dalam hal ini dapat dipilih funktor invers $T^{-1}$ yang memenuhi
	$T^{-1}T=\identity_{\mathcal D}=TT^{-1}$. Syarat tersebut mencakup semua
	situasi yang benar-benar ditangani dalam buku ini dan menghindari
	seluk-beluk potensial teori $2$-kategori, sehingga pembaca boleh saja
	membuat asumsi yang sama. Namun, kajian kategori model stabil atau kategori
	$\infty$ stabil tidak dapat menghindari kasus yang bukan isomorfisme,
	terutama dalam topologi.}
	Ekuivalensi $T$ disebut \emph{funktor pergeseran} dari $\mathcal{D}$.
	Selain itu, dipilih kuasi-invers $T^{-1}$ bagi $T$ beserta isomorfisme
	$T^{-1}T\simeq\identity_{\mathcal D}\simeq TT^{-1}$ sehingga data tersebut
	membentuk apa yang disebut ekuivalensi adjoin dalam
	\cite[Teorema~2.6.12]{Li1}.

% segment-id: o014.li2.chapter4.triangulated-category-definition.l001
	\begin{itemize}
% segment-id: o014.li2.chapter4.triangulated-category-definition.i001
		\item Suatu funktor dari $(\mathcal{D},T)$ ke
		$(\mathcal{D}',T')$ ialah funktor $F:\mathcal{D}\to\mathcal{D}'$
		beserta isomorfisme yang ditetapkan $FT\simeq T'F$. Menurut kebiasaan,
		isomorfisme itu tidak dituliskan dan data tersebut disingkat menjadi
		$F:(\mathcal{D},T)\to(\mathcal{D}',T')$.
% segment-id: o014.li2.chapter4.triangulated-category-definition.i002
		\item Suatu morfisme di antara dua funktor
% segment-id: o014.li2.chapter4.triangulated-category-definition.g001
		$\begin{tikzcd}
			(\mathcal{D},T) \arrow[r,shift left,"F"]
			\arrow[r,shift right,"{F'}"'] & (\mathcal{D}',T')
		\end{tikzcd}$
		ialah morfisme $\varphi:F\to F'$ yang membuat diagram berikut
		komutatif:
% segment-id: o014.li2.chapter4.triangulated-category-definition.g002
		\[\begin{tikzcd}
			FT \arrow[r,"\varphi T"] \arrow[d,"\sim"' sloped] &
			F'T \arrow[d,"\sim" sloped] \\
			T'F \arrow[r,"{T'\varphi}"'] & T'F'.
		\end{tikzcd}\]
% segment-id: o014.li2.chapter4.triangulated-category-definition.i003
		\item Suatu subkategori dari kategori dengan pergeseran
		$(\mathcal{D},T)$ ialah subkategori $\mathcal{D}'$ dari $\mathcal{D}$
		yang tertutup terhadap $T$ dan sedemikian sehingga
		$T':=T|_{\mathcal{D}'}$ menjadikan $(\mathcal{D}',T')$ kategori dengan
		pergeseran.
	\end{itemize}

% segment-id: o014.li2.chapter4.triangulated-category-definition.p002
	Jika $\mathcal{D}$ kategori aditif dan $T$ funktor aditif (hal ini
	berlaku secara otomatis; lihat Korolari~\ref{prop:automatic-additivity}),
	maka $(\mathcal{D},T)$ disebut \emph{kategori aditif dengan pergeseran};
	funktor di antara kategori-kategori tersebut juga disyaratkan aditif.
	Setelah suatu gelanggang komutatif $\Bbbk$ ditetapkan, definisi yang sama
	berlaku dalam situasi $\Bbbk$-linear.
\end{definisi}

% segment-id: o014.li2.chapter4.triangulated-category-definition.p003
Jika $(\mathcal{D},T)$ kategori dengan pergeseran, maka
$(\mathcal{D}^{\opp},T^{-1})$ juga demikian. Menurut kebiasaan, data
$(\mathcal{D},T)$ sering disingkat menjadi $\mathcal{D}$. Konsep berikut
sangat berguna.

% segment-id: o014.li2.chapter4.triangulated-category-definition.d002
\begin{definisi}[Morfisme berderajat]\label{def:morphism-with-degree}
	\index[sym1]{mdegree@$X \xrightarrow{+m} Y$}
	Suatu morfisme berderajat $m$ pada kategori dengan pergeseran
	$(\mathcal{D},T)$ ialah morfisme berbentuk $f:X\to T^mY$ dan juga
	dituliskan $f:X\xrightarrow{+m}Y$. Jika diberikan
	$X\xrightarrow[f]{+m}Y\xrightarrow[g]{+n}Z$, komposisi $gf$
	didefinisikan sebagai morfisme berderajat $m+n$, yaitu $T^m(g)f$.
	Diagram komutatif juga dapat dibicarakan untuk morfisme-morfisme jenis ini.

	Dalam situasi multivariabel, andaikan $\mathcal{D}$ dilengkapi
	automorfisme $T_1,\ldots,T_n$ yang berkomutasi secara ketat,
	$T_iT_j=T_jT_i$. Suatu morfisme berderajat $(m_1,\ldots,m_n)$
	didefinisikan sebagai morfisme berbentuk
	$X\to T_1^{m_1}\cdots T_n^{m_n}Y$. Morfisme ini dituliskan sebagai
	$X\xrightarrow{+(m_1,\ldots,m_n)}Y$.
\end{definisi}

% segment-id: o014.li2.chapter4.triangulated-category-definition.p004
Sebagai contoh, untuk suatu kategori $\mathcal{A}$, kategori objek bergradasi
$(\mathcal{A}^{\Z},T)$ dari Definisi~\ref{def:graded-obj} merupakan kategori
dengan pergeseran. Secara lebih umum, kita dapat meninjau kategori objek
bergradasi $\Z^n$, yaitu $\mathcal{A}^{\Z^n}$, beserta funktor pergeseran
$T_1,\ldots,T_n$ untuk setiap variabel, lalu membahas morfisme berderajat di
antara objek-objek bergradasi $\Z^n$. Namun, kasus multivariabel bukan pokok
bahasan bab ini.
\index{objek bergradasi}

% segment-id: o014.li2.chapter4.triangulated-category-definition.ex001
\begin{contoh}\label{eg:cplx-as-translation-cat}
	Kategori kompleks $\cate{C}(\mathcal{A})$ pada kategori aditif
	$\mathcal{A}$, bersama funktor pergeseran $T:X\mapsto X[1]$ dari
	Definisi~\ref{def:translation-functor}, merupakan kategori aditif dengan
	pergeseran. Kategori $\cate{K}(\mathcal{A})$ dari
	Definisi~\ref{def:cplx-homotopy-cat} mewarisi funktor pergeseran dari
	$\cate{C}(\mathcal{A})$, yang tetap dituliskan $T$. Funktor kanonik
	$F:\cate{C}(\mathcal{A})\to\cate{K}(\mathcal{A})$ jelas memenuhi
	$FT=TF$. Hal yang sama berlaku, untuk setiap
	$\star\in\{+,-,\bdd\}$, bagi kategori $\cate{C}^{\star}(\mathcal{A})$
	(atau $\cate{K}^{\star}(\mathcal{A})$) yang diperkenalkan dalam
	Definisi~\ref{def:cplx-cat-variant} (atau
	Definisi~\ref{def:cplx-homotopy-cat-variant}).
\end{contoh}

% segment-id: o014.li2.chapter4.triangulated-category-definition.p005
Dalam bab ini kita terutama meninjau kategori aditif dengan pergeseran.

% segment-id: o014.li2.chapter4.triangulated-category-definition.d003
\begin{definisi}
	\index{segitiga (triangle)}
	Misalkan $(\mathcal{D},T)$ kategori aditif dengan pergeseran. Suatu
	\emph{segitiga} di dalamnya ialah rangkaian morfisme pada $\mathcal{D}$
	berbentuk
% segment-id: o014.li2.chapter4.triangulated-category-definition.q001
	\[ X\xrightarrow{f}Y\xrightarrow{g}Z\xrightarrow{h}TX. \]
	Suatu morfisme di antara dua segitiga ialah diagram komutatif berbentuk
% segment-id: o014.li2.chapter4.triangulated-category-definition.g003
	\[\begin{tikzcd}
		X \arrow[r,"f"] \arrow[d,"\alpha"'] &
		Y \arrow[r,"g"] \arrow[d,"\beta"] &
		Z \arrow[r,"h"] \arrow[d,"\gamma"] &
		TX \arrow[d,"T\alpha"] \\
		X' \arrow[r,"{f'}"'] & Y' \arrow[r,"{g'}"'] &
		Z' \arrow[r,"{h'}"'] & TX'.
	\end{tikzcd}\]
	Kedua barisnya adalah segitiga yang diberikan. Morfisme tersebut juga
	disingkat menjadi $(X,Y,Z)\to(X',Y',Z')$. Dengan cara yang sama dapat
	didefinisikan isomorfisme di antara segitiga.

% segment-id: o014.li2.chapter4.triangulated-category-definition.p006
	Konsep serupa tetap tersedia jika kategori aditif diganti dengan kategori
	$\Bbbk$-linear yang lebih umum, dengan $\Bbbk$ sebarang gelanggang
	komutatif.
\end{definisi}

% segment-id: o014.li2.chapter4.triangulated-category-definition.p007
Jelas bahwa segitiga pada $(\mathcal{D},T)$ sama dengan segitiga pada
$(\mathcal{D}^{\opp},T^{-1})$.

% segment-id: o014.li2.chapter4.triangulated-category-definition.r001
\begin{catatan}\label{rem:triangle-signs}
	Tinjau segitiga
	$X\xrightarrow{f}Y\xrightarrow{g}Z\xrightarrow{h}TX$ pada kategori
	$\Bbbk$-linear dengan pergeseran $(\mathcal{D},T)$, serta
	$\epsilon,\zeta,\eta\in\Bbbk^\times$. Jika
	$\epsilon\zeta\eta=1$, segitiga semula isomorfik dengan
	$X\xrightarrow{\epsilon f}Y\xrightarrow{\zeta g}Z
	\xrightarrow{\eta h}TX$, sebagaimana ditunjukkan oleh diagram komutatif
% segment-id: o014.li2.chapter4.triangulated-category-definition.g004
	\[\begin{tikzcd}
		X \arrow[r,"f"] \arrow[d,"1"'] &
		Y \arrow[r,"g"] \arrow[d,"\epsilon"] &
		Z \arrow[r,"h"] \arrow[d,"\epsilon\zeta"] &
		TX \arrow[d,"\epsilon\zeta\eta"] \\
		X \arrow[r,"\epsilon f"'] & Y \arrow[r,"\zeta g"'] &
		Z \arrow[r,"\eta h"'] & TX.
	\end{tikzcd}\]
\end{catatan}

% segment-id: o014.li2.chapter4.triangulated-category-definition.p008
Menurut notasi Definisi~\ref{def:morphism-with-degree}, suatu segitiga juga
dapat dituliskan $X\to Y\to Z\xrightarrow{+1}$, atau secara lebih bergambar
sebagai
% segment-id: o014.li2.chapter4.triangulated-category-definition.g005
$\begin{tikzcd}[row sep=small,column sep=tiny]
	& Y \arrow[rd] & \\
	X \arrow[ru] & & Z \arrow[ll,"{+1}"]
\end{tikzcd}$.
Dengan demikian, segitiga dapat dibayangkan sebagai rantai yang berpilin naik,
sedangkan morfisme di antara segitiga menyerupai struktur asam
deoksiribonukleat: panah-panah $\alpha$, $\beta$, $\gamma$, dan seterusnya pada
morfisme tersebut dapat dianalogikan dengan ikatan hidrogen di antara basa.
\index{asam deoksiribonukleat (deoxyribonucleic acid)}

% segment-id: o014.li2.chapter4.triangulated-category-definition.d004
\begin{definisi}[\mbox{Kategori bertriangulasi} dan
\mbox{funktor bertriangulasi}]
\label{def:triangulated-cat}
	\index{kategori bertriangulasi (triangulated category)}
	\index{kategori prabertriangulasi (pretriangulated category)}
	\index{funktor!bertriangulasi (triangulated)}
	\index{segitiga!terbedakan (distinguished / exact)}
	Suatu \emph{kategori bertriangulasi} ialah kategori aditif dengan
	pergeseran $(\mathcal{D},T)$ yang dilengkapi suatu kelas segitiga
	$\mathcal{H}$, yang disebut \emph{segitiga terbedakan}, dan memenuhi
	aksioma-aksioma berikut.
% segment-id: o014.li2.chapter4.triangulated-category-definition.l002
	\begin{enumerate}[(TR1)]
		\setcounter{enumi}{-1}
% segment-id: o014.li2.chapter4.triangulated-category-definition.i004
		\item Setiap segitiga yang isomorfik dengan suatu segitiga terbedakan
		juga merupakan segitiga terbedakan.
% segment-id: o014.li2.chapter4.triangulated-category-definition.i005
		\item Untuk setiap $X\in\Obj(\mathcal{D})$, segitiga
		$X\xrightarrow{\identity_X}X\to0\to TX$ merupakan segitiga
		terbedakan.
% segment-id: o014.li2.chapter4.triangulated-category-definition.i006
		\item Setiap morfisme $f:X\to Y$ dapat diperluas menjadi suatu
		segitiga terbedakan berbentuk
		$X\xrightarrow{f}Y\to Z\to TX$, dengan $Z\in\Obj(\mathcal{D})$.
% segment-id: o014.li2.chapter4.triangulated-category-definition.i007
		\item Segitiga
		$X\xrightarrow{f}Y\xrightarrow{g}Z\xrightarrow{h}TX$ merupakan
		segitiga terbedakan jika dan hanya jika ``rotasi berlawanan arah jarum
		jam''-nya,
% segment-id: o014.li2.chapter4.triangulated-category-definition.q002
		\[ Y\xrightarrow{g}Z\xrightarrow{h}TX\xrightarrow{-Tf}TY, \]
		merupakan segitiga terbedakan. Pola tandanya juga boleh dipilih sebagai
		$(-++)$, $(+-+)$, atau $(---)$; lihat
		Catatan~\ref{rem:triangle-signs}.
% segment-id: o014.li2.chapter4.triangulated-category-definition.i008
		\item Andaikan bagian bergaris penuh pada diagram berikut komutatif,
% segment-id: o014.li2.chapter4.triangulated-category-definition.g006
		\[\begin{tikzcd}
			X \arrow[r] \arrow[d,"\alpha"'] &
			Y \arrow[r] \arrow[d,"\beta"] &
			Z \arrow[r] \arrow[dashed,d,"\gamma"] &
			TX \arrow[dashed,d,"T\alpha"] \\
			X' \arrow[r] & Y' \arrow[r] & Z' \arrow[r] & TX'.
		\end{tikzcd}\]
		Jika kedua barisnya segitiga terbedakan, terdapat morfisme
		$\gamma:Z\to Z'$ seperti ditunjukkan oleh garis putus-putus sehingga
		seluruh diagram yang memuat $\alpha$, $\beta$, $\gamma$, dan $T\alpha$
		komutatif. Dengan demikian diperoleh suatu morfisme segitiga.
% segment-id: o014.li2.chapter4.triangulated-category-definition.i009
		\item Andaikan diberikan segitiga-segitiga terbedakan
% segment-id: o014.li2.chapter4.triangulated-category-definition.q003
		\begin{gather*}
			X\xrightarrow{f}Y\xrightarrow{h}Z'\to TX,\\
			Y\xrightarrow{g}Z\xrightarrow{k}X'\to TY,\\
			X\xrightarrow{gf}Z\xrightarrow{m}Y'\to TX.
		\end{gather*}
		Terdapat segitiga terbedakan
% segment-id: o014.li2.chapter4.triangulated-category-definition.q004
		\[ Z'\xrightarrow{u}Y'\xrightarrow{v}X'\xrightarrow{w}TZ' \]
		sedemikian sehingga diagram berikut komutatif:
% segment-id: o014.li2.chapter4.triangulated-category-definition.g007
		\[\begin{tikzcd}
			X \arrow[r,"f"] \arrow[equal,d] &
			Y \arrow[r,"h"] \arrow[d,"g"] &
			Z' \arrow[r] \arrow[d,"u"] & TX \arrow[equal,d] \\
			X \arrow[r,"gf"] \arrow[d,"f"'] &
			Z \arrow[r,"m"] \arrow[equal,d] &
			Y' \arrow[r] \arrow[d,"v"] & TX \arrow[d,"Tf"] \\
			Y \arrow[r,"g"] \arrow[d,"h"'] &
			Z \arrow[r,"k"] \arrow[d,"m"] &
			X' \arrow[r] \arrow[equal,d] & TY \arrow[d,"Th"] \\
			Z' \arrow[r,"u"'] & Y' \arrow[r,"v"'] &
			X' \arrow[r,"w"'] & TZ'.
		\end{tikzcd}\]
	\end{enumerate}

% segment-id: o014.li2.chapter4.triangulated-category-definition.p009
	Data $(\mathcal{D},T,\mathcal{H})$ yang hanya memenuhi (TR0)--(TR4)
	disebut \emph{kategori prabertriangulasi}. Dalam notasi, data $T$ dan
	$\mathcal{H}$ sering tidak dituliskan.

% segment-id: o014.li2.chapter4.triangulated-category-definition.p010
	Suatu \emph{funktor bertriangulasi} dari $\mathcal{D}$ ke
	$\mathcal{D}'$ ialah funktor $F:\mathcal{D}\to\mathcal{D}'$ di antara
	kategori aditif dengan pergeseran yang memetakan segitiga terbedakan ke
	segitiga terbedakan. Morfisme di antara funktor bertriangulasi sama dengan
	morfisme dalam Definisi~\ref{def:category-with-translation}; dengan itu,
	kuasi-invers suatu funktor bertriangulasi dapat didefinisikan. Jika dua
	funktor bertriangulasi
% segment-id: o014.li2.chapter4.triangulated-category-definition.g008
	$\begin{tikzcd}
		\mathcal{D} \arrow[shift left,r,"F"] &
		\mathcal{D}' \arrow[shift left,l,"G"]
	\end{tikzcd}$
	saling kuasi-invers, keduanya disebut memberikan suatu \emph{ekuivalensi}
	di antara $\mathcal{D}$ dan $\mathcal{D}'$.
\end{definisi}

% segment-id: o014.li2.chapter4.triangulated-category-definition.p011
Jika $(\mathcal{D},T)$ bersifat $\Bbbk$-linear, semua definisi di atas
mempunyai perluasan yang bersesuaian.

% segment-id: o014.li2.chapter4.triangulated-category-definition.r002
\begin{catatan}\index{aksioma oktahedral (octahedral axiom)}
	Dengan notasi Definisi~\ref{def:morphism-with-degree}, aksioma (TR5) dapat
	dituliskan kembali sebagai diagram
% segment-id: o014.li2.chapter4.triangulated-category-definition.g009
	\[\begin{tikzcd}[column sep=large,row sep=large]
		& Y' \arrow[dashed,rd,"v"] \\
		Z' \arrow[dashed,ru,"u"] \arrow[d,"+1" description] &&
		X' \arrow[dashed,ll,"+1" description,"w" inner sep=0.6em]
		\arrow[ldd,"+1" description,near end] \\
		X \arrow[rd,"f"'] && Z \arrow[u,"k"'] \\
		& Y \arrow[ru,"g"'] \arrow[luu,"h"',near start] &
		\arrow[from=1-2,to=3-1,crossing over,"+1" description,near start]
		\arrow[from=3-1,to=3-3,crossing over]
		\arrow[from=3-3,to=1-2,crossing over,"m",near end]
	\end{tikzcd}\]
	Bagian bergaris putus-putus melambangkan morfisme yang keberadaannya
	dinyatakan oleh aksioma. Empat muka oktahedron ini merupakan segitiga
	terbedakan, sedangkan keempat muka lainnya komutatif; perinciannya
	diserahkan kepada pembaca. Karena itu, (TR5) juga disebut
	\emph{aksioma oktahedral}.
\end{catatan}

% [Yekutieli, Remark 5.6.12]
% segment-id: o014.li2.chapter4.triangulated-category-definition.r003
\begin{catatan}[Dualitas]\label{rem:triangulated-op}
	Misalkan $(\mathcal{D},T,\mathcal{H})$ kategori prabertriangulasi (atau
	kategori bertriangulasi). Pandang $T^{-1}$ sebagai funktor dari
	$\mathcal{D}^{\opp}$ ke dirinya sendiri dan tuliskan funktor itu sebagai
	$S$. Definisikan suatu kelas segitiga $\mathcal{K}$ pada
	$(\mathcal{D}^{\opp},S)$ sebagai berikut. Jika
% segment-id: o014.li2.chapter4.triangulated-category-definition.q005
	\[ [X\xrightarrow{f}Y\xrightarrow{g}Z\xrightarrow{h}TX]
		\;\in\mathcal{H}, \]
	maka rotasi searah jarum jamnya
	$[T^{-1}Z\xrightarrow{-T^{-1}h}X\xrightarrow{f}Y\xrightarrow{g}Z]
	\in\mathcal{H}$ dapat dipandang sebagai segitiga pada
	$\mathcal{D}^{\opp}$,
% segment-id: o014.li2.chapter4.triangulated-category-definition.q006
	\begin{equation}\label{eqn:op-triangle}
		Z\xrightarrow{g^{\opp}}Y\xrightarrow{f^{\opp}}X
		\xrightarrow{-(T^{-1}h)^{\opp}}SZ.
	\end{equation}
	Definisikan $\mathcal{K}$ sebagai kelas semua segitiga berbentuk
	\eqref{eqn:op-triangle}. Mudah diverifikasi bahwa
	$(\mathcal{D}^{\opp},S,\mathcal{K})$ menjadi kategori
	prabertriangulasi (atau kategori bertriangulasi). Juga jelas bahwa
	menjalankan prosedur ini dua kali mengembalikan
	$(\mathcal{D},T,\mathcal{H})$ semula.
\end{catatan}

% segment-id: o014.li2.chapter4.triangulated-category-definition.p012
Konsep subkategori prabertriangulasi dan subkategori bertriangulasi akan
diperkenalkan dalam
Definisi--Proposisi~\sourcecrossref{def:triangulated-subcat}{pada \S~4.2};
untuk sementara pembahasannya ditunda.
